% --------------------------------------------------------------
% This is all preamble stuff that you don't have to worry about.
% Head down to where it says "Start here"
% --------------------------------------------------------------
 
\documentclass[12pt]{article}
 
\usepackage[margin=0.8in]{geometry} % Margen
\usepackage{amsmath,amsthm,amssymb} % Mates
\usepackage{xcolor} % Colores

% Links color azul
\usepackage[
    colorlinks=true,
    urlcolor=blue
]{hyperref}
 
% Tablas
\usepackage{tabularx}
\usepackage{booktabs}
\usepackage{array}

 
\newenvironment{theorem}[2][Theorem]{\begin{trivlist}
\item[\hskip \labelsep {\bfseries #1}\hskip \labelsep {\bfseries #2.}]}{\end{trivlist}}
\newenvironment{lemma}[2][Lemma]{\begin{trivlist}
\item[\hskip \labelsep {\bfseries #1}\hskip \labelsep {\bfseries #2.}]}{\end{trivlist}}
\newenvironment{exercise}[2][Exercise]{\begin{trivlist}
\item[\hskip \labelsep {\bfseries #1}\hskip \labelsep {\bfseries #2.}]}{\end{trivlist}}
\newenvironment{reflection}[2][Reflection]{\begin{trivlist}
\item[\hskip \labelsep {\bfseries #1}\hskip \labelsep {\bfseries #2.}]}{\end{trivlist}}
\newenvironment{proposition}[2][Proposition]{\begin{trivlist}
\item[\hskip \labelsep {\bfseries #1}\hskip \labelsep {\bfseries #2.}]}{\end{trivlist}}
\newenvironment{corollary}[2][Corollary]{\begin{trivlist}
\item[\hskip \labelsep {\bfseries #1}\hskip \labelsep {\bfseries #2.}]}{\end{trivlist}}
 
\begin{document}
 
% --------------------------------------------------------------
%                         Start here
% --------------------------------------------------------------
 
%\renewcommand{\qedsymbol}{\filledbox}
 
\title{
    \vspace{-1.5cm}
    \large \textbf{Centro de Investigación y Docencia Económicas}\\
    \textbf{Macroeconomía II: Temario}\\
    \textbf{LECO - Otoño 2026}
}

\date{}




\maketitle
{
\begin{flushleft}
\vspace{-1.5cm}
 \textbf{Profesor:} Arturo López Pérez (\href{mailto:alopezperez.2605@gmail.com}{alopezperez.2605@gmail.com})

 \vspace{.1cm}
 \textbf{Laboratorista:} Jibran Elias Mejia (\href{mailto:jibran.elias@alumnos.cide.edu}{jibran.elias@alumnos.cide.edu})

 \vspace{.1cm}
 \textbf{Horario:} Martes y Jueves de 8:00 hrs a 9:30 hrs.
\end{flushleft}
}




\noindent Este es un curso de \textit{macroeconomía moderna} tal como se estudia en la investigación académica actual. El enfoque se basa en modelos de equilibrio general con fundamentos microeconómicos.
% , desde entornos de competencia perfecta hasta modelos con fricciones e imperfecciones de mercado.

\vspace{0.5cm}

\noindent \textbf{Objetivos:} Desarrollar las herramientas conceptuales y analíticas necesarias para construir, resolver e interpretar modelos macroeconómicos microfundamentados. Al finalizar el curso, el estudiante deberá ser capaz de identificar y comprender los supuestos y mecanismos que subyacen a los principales resultados teóricos, así como interpretar datos macroeconómicos y evidencia empírica a la luz de los modelos estudiados en clase.


\vspace{0.5cm}

\noindent \textbf{Prerequisitos:} El curso utiliza de manera intensiva herramientas de optimización estática con restricciones, particularmente el método de Lagrange, así como conceptos fundamentales de microeconomía. No se requiere conocimiento de optimización dinámica.

\vspace{0.5cm}

\noindent \textbf{Evaluación:}
\begin{itemize}
    \item Exámen Parcial 30\%
    \item Exámen Final 50\%
    \item Laboratorio 20\%
\end{itemize}

\vspace{0.5cm}

\noindent \textbf{Referencias:}
El curso estará guiado principalmente por los siguientes textos.
\begin{itemize}
    \item Williamson, Stephen D. (2018). \textit{Macroeconomics}. 6th Edition. Pearson Education Limited. (\textbf{WS})
    \item Kurlat, Pablo (2020). \textit{A Course in Modern Macroeconomics}. Self-published Lecture Notes. Stanford University. (\textbf{KP})

    \item Schmitt-Grohé, S., Uribe, M., \& Woodford, M. (2022). \textit{International Macroeconomics: A Modern Approach}. Princeton University Press. (\textbf{SUW})
\end{itemize}

La primera parte del curso se dedica a la formulación y el análisis de modelos macroeconómicos para una economía cerrada. La segunda parte estudia una economía abierta. En la primera parte, el curso sigue principalmente la exposición de WS, mientras que KP se utiliza para profundizar en los aspectos más técnicos de los modelos. Para la segunda parte, el curso sigue la exposición de SUW.

Aunque el curso tiene un contenido primordialmente teórico, el análisis de cada tema se complementará con evidencia empírica y aplicaciones provenientes de la literatura académica. Por ello, además de los libros de texto, una parte importante de las clases estará constituida por artículos de investigación que estudian empíricamente los mecanismos y evalúan las predicciones de los modelos examinados en clase.

\newpage

\subsection*{Syllabus}

\noindent El contenido del curso se distribuye en 29 sesiones de hora y media durante 15 semanas efectivas de clase, comenzando el 25 de agosto y concluyendo el 3 de diciembre. 

\vspace{0.1cm}


\subsubsection*{Primera Parte: Economía Cerrada}
 \begin{enumerate}
     \item Sesión de Bienvenida
     \begin{enumerate}
         \item[1.1] Introducción al curso
     \end{enumerate}
     \item La Conducta del Consumidor y de la Empresa (WS, Cap. 4)
     \begin{enumerate}
        \item[2.1] El consumidor representativo
        \item[2.2] Elección de ocio y trabajo
        \item[2.3] La empresa representativa
        \item[2.4] Maximización de beneficios
     \end{enumerate}
     \item Equilibrio General en un Periodo (WS, Cap. 5; KP, Cap. 9)
     \begin{enumerate}
         \item[3.1] Equilibrio competitivo
         \item[3.2] Eficiencia de Pareto 
         \item[3.3] Teoremas del bienestar 
     \end{enumerate}
     \item Elección Intertemporal de Consumo y Ahorro (KP, Cap. 6; WS, Cap. 9)
     \begin{enumerate}
         \item[4.1] Un modelo de dos periodos
         \begin{enumerate}
             \item[4.1.1] Efectos de la tasa de interés
             \item[4.1.2] La Hipótesis del Ingreso Permanente (PIH)
             \item[4.1.3] Impuestos y Equivalencia Ricardiana
             \item[4.1.4] Ahorro precautorio
         \end{enumerate}
         \item[4.2] Restricciones de liquidez
         \item [4.3] Teorías conductuales
         \begin{enumerate}
             \item[4.3.1] Sesgo al presente
             \item[4.3.2] Miopía
         \end{enumerate}
     \end{enumerate}
     \item Inversión y Acumulación de Capital (KP, Cap. 8; WS, Cap. 11)
     \begin{enumerate}
         \item[5.1] El producto marginal del capital y la inversión agregada
         \item[5.2] Un modelo de dos periodos
         \begin{enumerate}
             \item[5.2.1] Equilibrio parcial
             \item[5.2.2] Equilibrio general
         \end{enumerate}
     \end{enumerate}
     \item Equilibrio General Dinámico (KP, Cap. 9)
     \begin{enumerate}
         \item[6.1] Equilibrio general en una economía con dos periodos
         \item[6.2] Teoremas del bienestar
         \item[6.3] Equilibrio general en una economía con horizonte infinito
     \end{enumerate}
 \end{enumerate}

 
\begin{enumerate}
\setcounter{enumi}{6}
    \item Dinero e Inflación (WS, Cap. 12; KP, Cap. 10-11; DLS, Cap. 8)
    \begin{enumerate}
        \item[7.1] La oferta monetaria 
        \item[7.2] La demanda monetaria
        \item[7.3] Equilibrio en el mercado de dinero
        \item[7.4] La teoría cuantitativa del dinero
        \item[7.5] Un modelo monetario intertemporal: \textit{Cash-In-Advance}
    \end{enumerate}
    \item Modelo de Ciclos Reales de Negocios (KP, Cap. 12-13; DLS, Cap. 9; WS, Cap. 13)
    \begin{enumerate}
        \item[8.1] Hechos estilizados del ciclo económico
        \item[8.2] Breve repaso histórico: Keynesianismo, expectativas racionales y la crítica de Lucas
        \item[8.3] Un modelo de dos periodos
        \begin{enumerate}
            \item[8.3.1] Choques exógenos 
        \end{enumerate}
        % \item[8.4] Un modelo con horizonte infinito
        % \begin{enumerate}
        %     \item[8.4.1] Choques exógenos y funciones impulso-respuesta (IRF)
        % \end{enumerate}
    \end{enumerate}
    \item Modelo Nuevo-Keynesiano I (KP, Cap. 14; WS, Cap. 14)
    \begin{enumerate}
        \item[9.1] Competencia monopolística
        \item[9.2] Rigidez de precios
        \item[9.3] Un modelo de dos periodos
        \begin{enumerate}
            \item[9.3.1] Equilibrio general
        \item[9.3.2] Modelo IS-LM microfundamentado
        \item[9.3.3] Choques exógenos
        \item[9.3.4] Pérdidas de eficiencia
        \end{enumerate}
    \end{enumerate}
    \item Modelo Nuevo-Keynesiano II: Política Monetaria y Fiscal (KP, Cap. 15; WS, Cap. 15)
    \begin{enumerate}
        \item[10.1] Política fiscal
        \item[10.2] Política monetaria
        \begin{enumerate}
            \item[10.2.1] Inflación en el modelo Nuevo-Keynesiano
            \item[10.2.2] Curva de Phillips Aumentada
            \item[10.2.3] Expectativas racionales y $r^*$
        \end{enumerate}
        \item[10.3] Reglas de política monetaria
        \item[10.4] Trampa de liquidez
        \begin{enumerate}
            \item[10.4.1] La Cota Inferior Cero (ZLB)
            \item [10.4.2] Política monetaria no convencional
            \begin{enumerate}
                \item[10.4.2.1] Felixibilización cuantitativa (QE)
                \item[10.4.2.2] Guía prospectiva (\textit{Forward-Guidance})
            \end{enumerate}
        \end{enumerate}
    \end{enumerate}
\end{enumerate}

\newpage
\begin{enumerate}
\setcounter{enumi}{10}
    \item Desempleo (WS, Cap. 6; KP, Cap. 7)
     \begin{enumerate}
         \item[11.1] Modelo de Diamond-Mortensen-Pissarides
     \end{enumerate}
\end{enumerate}


\subsubsection*{Segunda Parte: Economía Abierta}


\begin{enumerate}
\setcounter{enumi}{11}
    \item Determinantes de la Cuenta Corriente (SUW, Cap. 2-3; OR, Cap. 1)
    \begin{enumerate}
    \item[12.1] Desequilibrios mundiales
    \item[12.2] Sostenibilidad de la cuenta corriente
    \item[12.3] Teoría intertemporal de la cuenta corriente
    \end{enumerate}
    \item Términos de Intercambio, la Tasa de Interés Mundial, Tarifas, y la Cuenta Corriente (SUW, Cap. 4)
    \begin{enumerate}
        \item[13.1] Choques en los términos de intercambio
        \item[13.2] Choques en la tasa de interés mundial
        \item[13.3] Tarifas en las importaciones
    \end{enumerate}
    \item Determinantes del Tipo de Cambio Real (SUW, Cap. 9-10; OR, Cap. 4)
    \begin{enumerate}
        \item[14.1] El tipo de cambio real y la paridad de poder de compra (PPP)
        \item[14.2] El modelo TNT
    \end{enumerate}
    \item Rigidez Nominal, Política Monetaria y Desempleo (SUW, Cap. 13; OR, Cap. 9)
    \begin{enumerate}
        \item[15.1] Dos tipos de rigidez nominal: precios domésticos y salarios nominales
        \item[15.2] El modelo TNT-DNWR
        \begin{enumerate}
            \item[15.2.1] Regímenes cambiarios
            \item[15.2.2] \textit{Overshooting}
        \end{enumerate}
        \item[15.3] La transmisión internacional de la política monetaria
        \begin{enumerate}
            \item[15.3.1] El \textit{trilemma} de la política monetaria
            \item[15.3.2] \textit{Dilemma not Trilemma}: El ciclo financiero global y la política monetaria
            \item[15.3.3] Coordinación internacional de política monetaria
        \end{enumerate}
    \end{enumerate}
\end{enumerate}

 % \subsubsection*{Tercera Parte: Temas Selectos de Macroeconomía\footnote{Condicional al avance del curso}}
 % \begin{enumerate}
 % \setcounter{enumi}{15}
 %    \item \textit{Twin-Deficits}: Déficits Fiscal y de Cuenta Corriente (SUW, Cap. 8)
 %    \begin{enumerate}
 %        \item[16.1] Una economía abierta con gobierno
 %        \item[16.2] Equivalencia Ricardiana
 %        \item[16.3] Gasto público y cuenta corriente
 %        \item[16.4] Optimalidad de los \textit{Twin-Deficits}
 %    \end{enumerate}
 %    \item Fricciones Financieras (WS, Cap. 10)
 %    \begin{enumerate}
 %        \item[17.1] Información asimétrica y crisis financieras
 %        \item[17.2] Precios de activos y burbujas financieras
 %    \end{enumerate}
 % \end{enumerate}







\newpage
\noindent\textbf{\large Calendario del curso}

\small
\vspace{0.15cm}

\noindent
\begin{tabularx}{\textwidth}{
    @{}
    >{\raggedright\arraybackslash}p{2.2cm}
    >{\raggedright\arraybackslash}p{4.1cm}
    >{\raggedright\arraybackslash}X
    @{}
}
\textit{Semana} & \textit{Sesiones} & \textit{Tema} \\
\midrule

Semana 1  & 25 y 27 de agosto       & Introducción al curso. La Conducta del Consumidor y de la Empresa (I). \\
Semana 2  & 1 y 3 de septiembre     & La Conducta del Consumidor y de la Empresa (II). Equilibrio General. \\
Semana 3  & 8 y 10 de septiembre    & Elección Intertemporal de Consumo y Ahorro. \\
Semana 4  & 15 y 17 de septiembre   & Inversión y Acumulación de Capital. Equilibrio General Dinámico. \\
Semana 5  & 22 y 24 de septiembre   & Modelo de Ciclos Reales de Negocios. \\
Semana 6  & 29 de sept. y 1 de oct. & Modelo Nuevo-Keynesiano. \\
\textbf{Semana 7}
& \textbf{6 y 8 de octubre}
& \textbf{Repaso (martes 6 de oct.). Examen parcial (jueves 8 de oct.)} \\
Semana 8  & 13 y 15 de octubre      & Desempleo. \\
Semana 9  & 20 y 22 de octubre      & Determinantes de la Cuenta Corriente (I). \\
Semana 10 & 27 y 29 de octubre      & Determinantes de la Cuenta Corriente (II). \\
Semana 11 & 3 y 5 de noviembre      & Términos de Intercambio, la Tasa de Interés Mundial, Tarifas y la Cuenta Corriente (I). \\
Semana 12 & 10 y 12 de noviembre    & Términos de Intercambio, la Tasa de Interés Mundial, Tarifas y la Cuenta Corriente (II). \\
Semana 13 & 17 y 19 de noviembre    & Determinantes del Tipo de Cambio Real. \\
Semana 14 & 24 y 26 de noviembre    & Rigidez Nominal, Política Monetaria y Desempleo (I). \\
Semana 15 & 1 y 3 de diciembre      & Rigidez Nominal, Política Monetaria y Desempleo (II). \\
          & \textbf{Del 7 al 18 de diciembre}
          & \textbf{Examen final (fecha por determinar)} \\

\bottomrule
\end{tabularx}

\vspace{0.1cm}

\noindent
\textit{\small Nota: La distribución de temas entre parciales puede cambiar dependiendo del avance del curso.}

 
\end{document}